\documentclass{article}
\usepackage[T1]{fontenc}
\usepackage{lmodern}
\usepackage{graphicx}
\usepackage{amsmath,amssymb}
\usepackage[a4paper,margin=2cm]{geometry}
\usepackage[absolute,overlay]{textpos}
\setlength{\TPHorizModule}{1mm}
\setlength{\TPVertModule}{1mm}
\usepackage[indonesian]{babel}
\usepackage{hyperref}
\setlength{\emergencystretch}{2em}
% unit: Halaman 33

\begin{document}

% === Halaman 33 (llm) ===

\begin{itemize}
    \item \textbf{Contoh:} Misalkan $P(2, 4)$ dan $Q(5, 9)$ adalah dua buah titik pada kurva eliptik $y^2 \equiv x^3 + x + 6 \pmod{11}$. Tentukan $P + Q$ dan $2P$.
\end{itemize}

\textbf{Jawab:}\
(a) $\color{red}{P + Q = R}$
\[
\begin{aligned}
\text{m} & \equiv (9 - 4)/(5 - 2) \pmod{11} = 5/3 \pmod{11} = 5 \cdot 3^{-1} \pmod{11} \\
& = 5 \cdot 4 \pmod{11} \equiv 9 \pmod{11}
\end{aligned}
\]
$P + Q = R$, koordinat Titik R:
\[
\begin{aligned}
x_r & \equiv \text{m}^2 - x_p - x_q \pmod{11} \equiv 81 - 2 - 5 \pmod{11} \equiv 8 \pmod{11} \\
y_r & \equiv \text{m}(x_p - x_r) - y_p \pmod{11} \equiv 9(2 - 8) - 4 \pmod{11} \equiv -58 \pmod{11} \\
& \equiv 8 \pmod{11}
\end{aligned}
\]

Jadi, $R(8, 8)$

\end{document}
